% !TeX program = xelatex
% UTF-8 transcription of all four source pages, in handwritten page order.
% Source wording and medical claims are retained without editorial correction.
% pic01.jpg = 20261011_184725377.jpg
% pic02.jpg = 20261011_184737585.jpg
% pic03.jpg = 20261011_184742562.jpg
% pic04.jpg = 20261011_184751731.jpg
% Compile directly with XeLaTeX, or input into a document using xepersian.
\ifdefined\settextfont
  \def\InfluenzaStandalone{0}
\else
  \def\InfluenzaStandalone{1}
  \documentclass[12pt,a4paper]{article}
  \usepackage[margin=2cm]{geometry}
  \usepackage{amsmath,amssymb}
  \usepackage{xepersian}
  \IfFontExistsTF{Amiri}{\settextfont{Amiri}}{\settextfont{FreeSerif}}
  \setlatintextfont{FreeSerif}
\fi
\providecommand{\en}[1]{\lr{#1}}
\newcommand{\influpage}[1]{\clearpage\section*{صفحهٔ #1}%
  {\footnotesize\en{No \ldots\quad Date \ldots}\par}}
\ifnum\InfluenzaStandalone=1 \begin{document}\fi

\influpage{۱}
\section*{آنفولانزا}
\subsection*{درمان آنفولانزا}
\begin{itemize}
\item در گذشته از مهارکننده‌های کانال \en{M-2} برای درمان آنفولانزا \en{A} استفاده می‌شد.
  \begin{itemize}
  \item آمانتادین، ریمانتادین.
  \item در حال حاضر مقاومت گسترده‌ای نسبت به این داروها وجود دارد.
  \end{itemize}
\item مهارکننده‌های نورآمینیداز: در حال حاضر درمان اصلی هستند.
  \begin{itemize}
  \item برای آنفولانزا \en{A} و \en{B}.
  \item بیشترین تأثیر در افرادی که آنفولانزا آن‌ها تازه شروع شده و با تست تأیید شده است.
  \item اگر ظرف \en{48 h} اول بیماری داده شوند علائم را از بین می‌برند.
  \item اندیکاسیون به عنوان پروفیلاکسی هم دارند / بیماری ۱--۲ روز زودتر رفع.
  \item همهٔ این داروها خطر مقاومت دارند.
    \begin{itemize}
    \item خصوصاً اگر طولانی‌مدت مصرف شوند: افراد نقص ایمنی.
    \end{itemize}
  \end{itemize}
\item گروه‌های هدف برای درمان:
  \begin{itemize}
  \item کودکان زیر ۲ سال.
  \item افراد بالای ۶۵ سال.
  \item افراد با بیماری‌های مزمن.
  \item بیماران سرکوب ایمنی.
  \item زنان باردار، زایمان ظرف ۲ هفتهٔ اخیر.
  \item بیماران زیر ۱۸ سال که آسپرین طولانی‌مدت دریافت می‌کنند.
  \item افراد با چاقی مرضی.
  \item افراد در مراکز نگهداری.
  \end{itemize}
\item اوسلتامیویر (خوراکی)، زانامیویر (اسپری بینی)، پرامیویر (وریدی).
  \begin{itemize}
  \item طول دورهٔ درمان: ۵ روز، روزی ۲ بار.
  \item اوسلتامیویر: انتخابی برای زنان باردار، از ۱ سالگی می‌توانیم بدهیم.
  \item زانامیویر در آسم و \en{COPD} کنتراندیکاسیون دارد.
    \begin{itemize}
    \item بالای ۷ سال \en{OK}، می‌توانیم بدهیم.
    \end{itemize}
  \end{itemize}
\item در بیمار بستری بلافاصله پس از تشخیص اوسلتامیویر را شروع می‌کنیم.
  \begin{itemize}
  \item اگر امکان خوراکی نباشد تک‌دوز پرامیویر می‌دهیم: بالای ۲ سال می‌توانیم بدهیم.
  \end{itemize}
\item بالوکساویر (\en{Baloxavir}): پروفیلاکسی، درمان آنفولانزا بدون عارضه.
  \begin{itemize}
  \item سن حداقل ۱۲ سال.
  \item باید ظرف ۴۸ ساعت نخست بیماری آغاز شود.
  \item تک‌دوز است.
  \item مؤثر برای آنفولانزا \en{A} و \en{B}.
  \end{itemize}
\end{itemize}
{\footnotesize\en{PAPCO}\quad ۱\par}

\influpage{۲}
% The first two bullets continue the Baloxavir entry on source page 1.
\begin{itemize}
\item استفاده در زنان باردار و شیرده توصیه نمی‌شود.
\item استفاده در گروه‌های پرخطر توصیه نمی‌شود: چون دیتا نداریم.
\end{itemize}
کنترل تب با \en{NSAID} جزئی از درمان بیماری است.

\subsection*{ایمن‌سازی علیه آنفولانزا}
\begin{itemize}
\item واکسیناسیون بهترین رویکرد برای پیشگیری از آنفولانزا است.
  \begin{itemize}
  \item در آمریکا در ماه آگوست و سپتامبر موجود می‌شود.
  \item شدت بیماری را کاهش می‌دهد.
  \end{itemize}
\item اندیکاسیون تزریق روتین سالانه:
  \begin{itemize}
  \item ۶ ماه به بالا.
  \item برای ۶۵ سال به بالا نسخهٔ \en{high-dose}.
  \end{itemize}
\item کودکان زیر ۹ سال اگر تاکنون حداقل ۲ بار واکسن نزده باشند باید ۲ نوبت تزریق کنند.
\item سابقهٔ حساسیت و آنافیلاکسی به هر یک از اجزاء واکسن کنتراندیکاسیون تزریق واکسن است.
\item وجود سابقهٔ سندرم گیلن‌باره ظرف ۶ هفته پس از واکسن آنفولانزا: احتیاط.
\item حساسیت به تخم‌مرغ: نیاز به احتیاط بیشتری ایجاد نمی‌کند.
\end{itemize}

\subsection*{تشخیص‌های افتراقی آنفولانزا}
\begin{itemize}
\item بیماری‌های تب‌دار تنفسی در \en{ddx} آنفولانزا قرار دارند.
\item کووید: از دست رفتن حاد حس بویایی و چشایی به نفع کووید است.
\end{itemize}

\subsection*{یافته‌های آزمایشگاهی}
\begin{itemize}
\item لکوپنی خفیف در آنفولانزا مورد انتظار است.
\item \en{$\mathrm{WBC}>15$} مطرح‌کنندهٔ سوار شدن عفونت باکتریال است.
% The source writes WBC > 15 without a unit or a multiplier.
\item تست‌های سرولوژیک ۲ هفته طول می‌کشند تا مثبت شوند.
\item اندیکاسیون \en{Rapid Test}، \en{NAAT}:
  \begin{itemize}
  \item همهٔ بیماران بستری، \en{high risk}.
  \item بیماران در مراکز نگهداری.
  \item هر بیماری که نتیجهٔ تست در شیوهٔ درمان ما مؤثر باشد.
  \end{itemize}
\end{itemize}

\subsection*{تظاهرات بالینی آنفولانزا}
\begin{itemize}
\item سرفه، گلودرد، رینوره، احتقان بینی.
\item شروع ناگهانی علائم (\en{Sudden Onset}).
\item شدت تب، لرز، ضعف، خستگی، بدن‌درد: بارزتر از سایر عفونت‌های ویروسی / مثل کووید.
\end{itemize}
{\footnotesize\en{PAPCO}\quad ۲\par}

\influpage{۳}
\begin{itemize}
\item علائم معمولاً ظرف ۴۸--۷۲ ساعت از مواجهه شروع می‌شوند.
\item طول بیماری ۲--۵ روز است ولی علائم تنفسی خصوصاً سرفه بیشتر طول می‌کشند.
\item یک تا چند زمانی پس از دورهٔ عفونی وجود دارد تا فرد به سطح عملکردی اولیه برگردد.
% The tightly written phrase «یک تا چند زمانی» is an uncertain reading.
% It is transcribed as written rather than expanded to a presumed duration.
\item در افراد مسن علائم تنفسی ممکن است خفیف‌تر باشند.
  \begin{itemize}
  \item ضعف و بی‌اشتهایی شدیدتر هستند.
  \end{itemize}
\item تعریق، لنفادنوپاتی غیر فوکوس، اریتم منتشر فارنکس: شایع هستند.
\item معاینهٔ ریه:
  \begin{itemize}
  \item رال منتشر، رونکای، ویز.
  \item یافته‌های لوکالیزه معمولاً بیانگر عفونت باکتریال ثانویه است.
  \end{itemize}
\item در آنفولانزا راش نداریم.
\item علائم گوارشی در آنفولانزا نادر است.
\end{itemize}

\subsection*{آنفولانزا عارضه‌دار}
\begin{itemize}
\item بیشتر در بالای ۶۵ سال، کودکان کم‌سن، بیماری زمینهٔ قلبی‌ریوی، سرکوب ایمنی، زنان در سه‌ماههٔ دوم و سوم بارداری.
\end{itemize}

\subsection*{عوارض ریوی}
\begin{itemize}
\item پنومونی:
  \begin{itemize}
  \item یافته‌های لوکالیزه در ریه.
  \item شواهد رادیوگرافیک پنومونی.
  \item افزایش دیس‌پنه، تب بالا، سیانوز.
  \item در نوع باکتریال استاف طلایی و پنوموکوک شایع هستند.
  \item در بیماری زمینه‌ای ریه (\en{COPD}) عفونت ثانویهٔ باکتریال شایع‌تر است.
  \end{itemize}
\item برخی سوش‌ها لارنگوتراکئوبرونشیت ایجاد می‌کنند، برونشیولیت، \en{Croup}.
  \begin{itemize}
  \item خصوصاً در کودکان.
  \end{itemize}
\item اوتیت مدیا ویروسی و باکتریایی خصوصاً در کودکان.
\end{itemize}

\subsection*{عوارض خارج ریوی}
\begin{itemize}
\item میوزیت شایع‌ترین درگیری خارج ریوی در آنفولانزا است.
  \begin{itemize}
  \item در آنفولانزا \en{B} شایع‌تر است.
  \item درد شدید عضلانی، افزایش \en{CPK}، میوگلوبینوری، تندرنس عضله.
  \end{itemize}
\item میوکاردیت / پریکاردیت: کمتر شایع هستند.
\item انسفالیت / انسفالوپاتی / میلیت عرضی (\en{TM}) / پارکینسون.
\item \en{Postinfectious Acute Demyelinating Encephalomyelitis}.
\item درگیری‌های نورولوژیک در کودکان شایع‌تر است.
  \begin{itemize}
  \item کودکان به طور شایع دچار تب و تشنج می‌شوند.
  \item گیلن‌باره.
  \end{itemize}
\end{itemize}
{\footnotesize\en{PAPCO}\quad ۳\par}

\influpage{۴}
\subsection*{گروه‌های پرریسک}
\begin{itemize}
\item کودکان ۶--۵۹ ماهه.
\item سن بالای ۵۰ سال.
\item زنان باردار.
\item بیماری زمینه‌ای.
\item بیماران نقص ایمنی.
\item \en{$\mathrm{BMI}>40$}.
\item نقص ایمنی به هر علتی.
\end{itemize}
طغیان آنفولانزا معمولاً در ابتدای زمستان شروع می‌شود و ۴--۵ هفته ادامه دارد.
{\footnotesize\en{PAPCO}\quad ۴\par}

\ifnum\InfluenzaStandalone=1
  \let\InfluenzaEnd\enddocument
\else
  \let\InfluenzaEnd\relax
\fi
\InfluenzaEnd
